\chapter{Introduction}
\label{ch:introduction}

\section{Context}

The telecommunications industry is undergoing a fundamental transformation driven by the convergence of 5G networks, cloud computing, and artificial intelligence. Network operators face increasing complexity in managing thousands of cell sites, millions of subscribers, and petabytes of telemetry data generated every day. Traditional reactive network operations---where engineers manually investigate alarms after service degradation occurs---are no longer sufficient to meet the stringent quality expectations of modern digital services.

Huawei Technologies, as a global leader in telecommunications infrastructure, has pioneered the concept of \gls{adn}, a framework that envisions progressively autonomous network operations across five levels (L1--L5), inspired by the autonomous driving paradigm in the automotive industry. At the heart of this vision lies the intelligent convergence of \gls{oss} and \gls{bss} data, enabling operators to correlate network performance with business impact in real time.

This project was conducted within Huawei Tunisia's Cloud IT / Sales-Solution department, with a focus on building an intelligence layer that bridges \gls{cem} (powered by Huawei SmartCare) and \gls{cvm} for Tunisie Telecom's network operations.

\section{Problem Statement}

Current telecommunications operations suffer from three critical limitations:

\begin{enumerate}
    \item \textbf{Reactive SLA Management:} \gls{sla} breaches are detected after they occur, resulting in penalty costs and subscriber churn. There is no predictive mechanism to anticipate breaches before they happen.

    \item \textbf{Siloed OSS-BSS Data:} Network performance metrics (throughput, latency, packet loss) and business metrics (revenue, data usage, churn risk) are analyzed in isolation. The correlation between network degradation and revenue impact remains invisible to operators.

    \item \textbf{Manual Anomaly Investigation:} Network anomalies and revenue irregularities (such as SIM box fraud) require manual investigation, consuming engineering resources and increasing mean time to resolution (MTTR).
\end{enumerate}

\section{Proposed Solution}

This project delivers a \textbf{Cloud-Native AI Operations Agent} that addresses these limitations through:

\begin{itemize}
    \item A \textbf{GradientBoosting regression model} that predicts SLA breach risk from aggregated network \gls{kpi} features, enabling proactive capacity management.
    \item Two \textbf{IsolationForest anomaly detection models} that identify degraded network conditions and suspicious revenue patterns in an unsupervised manner.
    \item A \textbf{Pearson--Spearman correlation engine} that quantifies OSS--BSS relationships, revealing how network quality impacts business metrics.
    \item An \textbf{L4 Autonomous Agent} with an integrated local \gls{llm} (Qwen2.5) that provides conversational intelligence, autonomous remediation recommendations, and closed-loop decision support.
    \item A \textbf{cloud-native microservices architecture} with seven containerized services, designed for direct portability to \gls{hcs}.
\end{itemize}

\section{Objectives}

The project pursues the following objectives:

\begin{enumerate}
    \item Design and implement a production-grade data pipeline that ingests, processes, and curates OSS and BSS data in 2-minute cycles.
    \item Train and deploy three \gls{ml} models for SLA risk prediction and anomaly detection.
    \item Build a real-time dashboard (Next.js 14) with advanced visualizations for all intelligence outputs.
    \item Implement an L4 Autonomous Agent aligned with the TM Forum \gls{adn} framework.
    \item Ensure cloud portability through Docker containerization and \gls{hcs} service mapping (OBS, RDS, ECS).
    \item Establish a \gls{cicd} pipeline with automated linting, testing, building, and security scanning.
\end{enumerate}

\section{Report Structure}

This report is organized as follows:

\begin{description}
    \item[Chapter 2] reviews the state of the art in autonomous network operations, CEM/CVM convergence, and relevant ML techniques.
    \item[Chapter 3] presents the requirements analysis and specification.
    \item[Chapter 4] details the system architecture and technology choices.
    \item[Chapter 5] describes the implementation of core services and the data pipeline.
    \item[Chapter 6] focuses on ML model training, evaluation, and deployment.
    \item[Chapter 7] presents the L4 Autonomous Agent with LLM integration.
    \item[Chapter 8] covers deployment, CI/CD, and Huawei Cloud Stack mapping.
    \item[Chapter 9] presents results, testing, and evaluation.
    \item[Chapter 10] concludes the report with a summary and future work.
\end{description}
